\قسمت{چرا \متن‌لاتین{HTTP2}؟}

در نسخه ۱.۱ از پروتکل \متن‌لاتین{HTTP} بحث \متن‌لاتین{Pipelining} مطرح شد، چه چیزی باعث شد که عملا این بحث بهبود ناچیزی در کارایی این پروتکل داشته باشد؟
در \متن‌لاتین{HTTP2}، مساله‌ی \متن‌لاتین{Pipelining} چگونه حل شده است؟

\نمره{۲}

\begin{پاسخ}

در \متن‌لاتین{Pipelining} در \متن‌لاتین{HTTP} نسخه ۱.۱ نیاز دارد که ترتیب درخواست‌ها حفظ شود
چرا که از روی خود بسته‌ها نمی‌توان ترتیب آن‌ها را تشخیص داد.
بنابراین مشکل بلاک شدن درخواست‌ها به دلیل پردازش کند یک درخواست
وجود دارد (ذکر این نکته به عنوان دلیل عدم استفاده از \متن‌لاتین{Pipelining} نیمی از نمره را دارد). این مساله در \متن‌لاتین{HTTP2} با استفاده از \متن‌لاتین{Multiplexing} حل شده است.
نیازی به توضیح استفاده از شناسه روی درخواست و تقاضا برای تشخیص ارتباط‌شان با یکدیگر نیست.

\end{پاسخ}
